\begin{table}[tbp]
  \centering
  \caption{Matriz de transición de pobreza, pobreza monetaria y
  multidimensional (IPM), 2010$\to$2013 y 2013$\to$2016.}
  \label{tab:matriz_transicion}
  \footnotesize
  \setlength{\tabcolsep}{4pt}
  \begin{tabular}{l rrr rrr}
    \toprule
    & \multicolumn{3}{c}{\textbf{Monetaria}} & \multicolumn{3}{c}{\textbf{IPM}} \\
    \cmidrule(lr){2-4}\cmidrule(lr){5-7}
    \textbf{Primera encuesta} & \textbf{No pobre} & \textbf{Pobre} & \textbf{Total} & \textbf{No pobre} & \textbf{Pobre} & \textbf{Total} \\
    \midrule
    \multicolumn{7}{l}{\textit{2010$\to$2013}} \\
    \quad No pobre & 2{,}366 (76.6\%) & 723 (23.4\%) & 3{,}089 & 5{,}690 (87.9\%) & 784 (12.1\%) & 6{,}474 \\
    \quad Pobre & 1{,}525 (29.7\%) & 3{,}604 (70.3\%) & 5{,}129 & 951 (54.5\%) & 793 (45.5\%) & 1{,}744 \\
    \addlinespace
    \multicolumn{7}{l}{\textit{2013$\to$2016}} \\
    \quad No pobre & 2{,}556 (80.1\%) & 635 (19.9\%) & 3{,}191 & 5{,}029 (90.3\%) & 542 (9.7\%) & 5{,}571 \\
    \quad Pobre & 1{,}422 (38.2\%) & 2{,}298 (61.8\%) & 3{,}720 & 701 (52.3\%) & 639 (47.7\%) & 1{,}340 \\
    \bottomrule
  \end{tabular}
  \begin{minipage}{0.95\textwidth}
    \vspace{4pt}
    \footnotesize \textit{Nota:} Número de hogares y, entre paréntesis,
    porcentaje de la fila: de los hogares en cada estado en la primera
    encuesta del periodo, qué porcentaje está en cada estado en la
    siguiente. Panel emparejado 1 a 1, con los mismos hogares para las dos
    definiciones de pobreza ($n=8{,}218$ en 2010$\to$2013 y $n=6{,}911$
    en 2013$\to$2016); excluye a los hogares que se dividieron entre olas
    (511 en 2010$\to$2013 y 1{,}190 en 2013$\to$2016). Sin ponderar.
    Fuente: cálculos propios.
  \end{minipage}
\end{table}